% !PRINTED 585
\weberpage{585}

Es ist dann nach (2)
\weq{8}{\begin{aligned}
\mu &= y + bx + 2a\omega x, \\
\mu\omega &= -2cx + (y-bx)\omega, \\
m &= (y+bx)(y-bx) + 4acx^2,
\end{aligned}}
woraus folgt, da\ss{}
\[
y + bx \text{ und } y - bx
\]
ungerade sind.

Daraus ergibt sich nach (4):
\weq{9}{\begin{aligned}
\boldsymbol{s}[\mu(u+1)] &= -\boldsymbol{s}(\mu u), \\
\boldsymbol{s}[\mu(u+\omega)] &= \boldsymbol{s}(\mu u).
\end{aligned}}
Die Funktion $\boldsymbol{s}(\mu u)$ hat also dieselben Perioden wie $\boldsymbol{s}(u)$,
und da $\boldsymbol{s}(\mu u)$ eine ungerade Funktion von $u$ ist, so kann sie rational durch $\boldsymbol{s}(u)$ ausgedrückt werden. Wir bezeichnen mit $A(\boldsymbol{s})$, $D(\boldsymbol{s})$ ganze rationale Funktionen von $\boldsymbol{s} = \boldsymbol{s}(u)$, die übrigens nur die geraden Potenzen von $\boldsymbol{s}$ enthalten, und setzen:
\weq{10}{\boldsymbol{s}(\mu u) = \mu\,\boldsymbol{s}(u)\,\dfrac{A(\boldsymbol{s})}{D(\boldsymbol{s})}.}
Die Werte von $\boldsymbol{s}(u)$, für die $\boldsymbol{s}(\mu u)$ Null oder unendlich wird, sind nach 3. algebraische Zahlen, und da sich aus diesen die Koeffizienten von $A$ und $D$ zusammensetzen lassen, so sind diese auch algebraische Zahlen. Um ihre Natur näher zu bestimmen, denken wir uns den Bruch (10) zunächst so erweitert, da\ss{} $D$ rationale Koeffizienten erhält. Das geschieht dadurch, da\ss{} wir Zähler und Nenner mit dem Produkt aller zum Nenner konjugierten Faktoren multiplizieren. Wir können also $A(\boldsymbol{s})$ und $D(\boldsymbol{s})$ in die Form setzen:
\weq{11}{\begin{aligned}
A(\boldsymbol{s}) &= A_1 + A_3\boldsymbol{s}^2 + A_5\boldsymbol{s}^4 + \cdots, \\
D(\boldsymbol{s}) &= D_1 + D_3\boldsymbol{s}^2 + D_5\boldsymbol{s}^4 + \cdots,
\end{aligned}}
worin die $D_1, D_3, \dots$, nach unserer Voraussetzung rationale Zahlen sind, während $A_1, A_3, \dots$, zu bestimmen sind. Setzen wir wie in \S\,155, (4), (6):
\[
w = \sqrt{\varkappa}\,v = 2\sqrt{\varkappa}\,Ku,
\]
so geht $\boldsymbol{s}(u)$ in $\sqrt{\varkappa}\ \mathrm{sn}\left(\dfrac{w}{\sqrt{\varkappa}}\right)$ über, und die Gleichung (10) ergibt:
\weq{12}{\dfrac{D\left[\sqrt{\varkappa}\ \mathrm{sn}\left(\dfrac{\omega}{\sqrt{\varkappa}}\right)\right]\ \sqrt{\varkappa}\ \mathrm{sn}\left(\dfrac{\mu w}{\sqrt{\varkappa}}\right)}{\mu\ \mathrm{sn}\left(\dfrac{w}{\sqrt{\varkappa}}\right)} = A_1 + A_3\boldsymbol{s}^2 + A_5\boldsymbol{s}^4 + \cdots}
% [source sic: D's argument prints \omega while the other arguments use w, and the denominator lacks the otherwise matching \sqrt{\varkappa}]
